The comparison between problems' difficulties is done by \textbf{reductions}, something that we have already seen in the first chapters. \vspace{7pt}

We would like to define a reduction between two decision problems (\textit{namely languages}) through a function $f$. Let's consider the following image.
\begin{center}
    % \includegraphics[width=0.5\textwidth]{img/img1_.png}
\end{center}
The function $f$ takes in input a point from the left-most set and maps it to a point in the right-most set: considering a starting point inside the language 
$L$, the function $f$ will land it inside the language $H$. Automatically we can say that $L$ belongs to the same complexity of $H$.
\begin{definition}
    The language $L$ is said to be \textbf{polynomial-time reducible} to another language $H$ if and only if there is a polytime computable function $f: \{0,1\}^*
    \rightarrow \{0,1\}^*$ such that $x \in L$ if and only if $f(x) \in H$.
\end{definition}
If we have the capabilities to define such relationship $f$ between the decision problems considered, we write $L \leq_p H$, meaning that: $L$ cannot be 
harder than $H$.